\uuid{3fY4}
\titre{Changements de variables}
\auteur{hanani ; mzari}
\organisation{exo7}
\licence{CC BY-NC-SA 4.0}
\datecreate{2020-01-20}
\niveau{L1}
\module{Analyse}
\chapitre{Calcul d'intégrales}
\sousChapitre{Primitives diverses}
\difficulte{3}
\theme{Primitives, changements de variables}
\exo7qcmid{lille-2.primitives.050}

\contenu{
\begin{question}
\qclassification{127.11}
\qtags{motcle=changements de variables}
Soient $a$ et $b$ deux réels et $f$ la fonction telle que $\displaystyle f(x)=\frac{1}{\sqrt{ax^2+b}}$. Parmi les affirmations suivantes, cocher celles qui sont vraies :
\begin{answers}
\bad{Si $a$ et $b$ sont strictement positifs, il existe une constante $C$ telle que $F(x)=C\arcsin \left(x\sqrt{a/b}\right)$ soit une primitive de $f(x)$.}
\good{Si $a$ et $b$ sont strictement positifs, il existe une constante $C$ telle que $F(x)=C\sqrt{ax^2+b}$ soit une primitive de $xf(x)$.}
\good{Si $a$ est strictement négatif et $b$ est strictement positif, il existe une constante $C$ telle que $F(x)=C\arcsin \left(x\sqrt{-a/b}\right)$ soit une primitive de $f(x)$.}
\bad{Si $a$ est strictement positif et $b$ est strictement négatif, il existe une constante $C$ telle que $F(x)=C\arcsin \left[x\sqrt{a/(-b)}\right]$ soit une primitive de $f(x)$.}
\end{answers}
\vskip2mm
\begin{explanations}
Il suffit de dériver : $\displaystyle \left[C\arcsin \left(x\sqrt{a/b}\right)\right]'=\frac{C\sqrt{a}}{\sqrt{b-ax^2}}\neq f(x)$ pour tout $C\in \Rr$. Par contre,
$$\left(C\sqrt{ax^2+b}\right)'=\frac{Cax}{\sqrt{ax^2+b}}=xf(x)\quad \mbox{si}\quad C=a^{-1}.$$
Si $a$ est strictement négatif et $b$ est strictement positif,
$$\left[C\arcsin \left(x\sqrt{-a/b}\right)\right]'=\frac{C\sqrt{-a}}{\sqrt{b+ax^2}}=f(x)\quad \mbox{si}\quad C\sqrt{-a}=1.$$
Par contre, si $a$ est strictement positif et $b$ est strictement négatif alors, pour tout $C\in \Rr$, $\left[C\arcsin \left(x\sqrt{a/(-b)}\right)\right]'\neq f(x)$.
\end{explanations}
\end{question}
}
